\documentclass[11pt,letterpaper]{article}
\usepackage[T1]{fontenc}
\usepackage[utf8]{inputenc}
\usepackage{lmodern}
\usepackage{amsmath,amssymb,amsthm}
\usepackage{microtype}
\usepackage[margin=0.9in]{geometry}
\usepackage{enumitem}
\usepackage{xcolor}
\usepackage{hyperref}
\usepackage{booktabs,longtable,array}
\hypersetup{colorlinks=true,linkcolor=black,urlcolor=blue!45!black,pdftitle={Late coefficients of three letter abelian squares},pdfauthor={dot},pdfsubject={A274600 proof, Stokes jump, and inverse sectors}}
\setlength{\parindent}{0pt}
\setlength{\parskip}{5pt plus 1pt}
\setlength{\emergencystretch}{2em}
\setlist{itemsep=3pt,topsep=5pt}
\providecommand{\tightlist}{\setlength{\itemsep}{0pt}\setlength{\parskip}{0pt}}
\allowdisplaybreaks[2]
% ed. (2026-10-02): unnumbered environment for ProveIt editorial notes (no counter is used; upright body).
\theoremstyle{remark}
\newtheorem*{ednote}{Editorial note (ProveIt, 2026-10-02)}
\title{\textbf{Late coefficients of three letter\\abelian squares}\\[0.5em]\large All fixed late orders and exact inverse Stokes sectors}
\author{A274600 research note}
\date{2 October 2026}
\begin{document}
\maketitle
\begin{abstract}
The coefficients in the inverse-power expansion of the three-letter abelian-square counts grow factorially. This note proves the explicitly labeled OEIS A274600 asymptotic conjecture and gives an effective expansion to every fixed late order. The proof uses the classical three-step random-walk density to identify the exact Borel transform, proves its required complex continuation, and extracts the logarithmic singularity by a keyhole contour. It also establishes the exact physical median and Stokes jump, Lambert-function index inversions, and a convergent generator of all fixed inverse exponential sectors. Exact-arithmetic replay and numerical checks accompany the proof. Novelty is stated only within the bounds of the documented source search.
\end{abstract}
\section{Target, notation, and principal
result}\label{target-notation-and-principal-result}

Write

\[
M_N=\sum_{a+b+c=N}\binom{N}{a,b,c}^{\!2}
    =\sum_{k=0}^N\binom Nk^2\binom{2k}k\qquad(\text{A002893}).
\]

Let \(a_n\) be the uniquely determined Poincaré coefficients in

\[
M_N\sim \frac{C_0 9^N}{N}\sum_{n\ge0}\frac{a_n}{(4N)^n},
\qquad C_0=\frac{3\sqrt3}{4\pi}.
\tag{1}
\]

These are A274600: \(1,-1,1,2,7,59,616,6992,\ldots\). The currently
inspected OEIS entry, authored by Václav Kotešovec on 10 November 2016,
explicitly conjectures

\[
a_n\sim\frac1{\pi\sqrt3}\left(\frac2{\log3}\right)^n(n-1)!.
\tag{2}
\]

This concerns the late coefficients of (1), not the leading asymptotic
of \(M_N\).

Put \(L=\log9\), \(\kappa=4/L=2/\log3\), \(C=1/(\pi\sqrt3)\), and define

\[
\begin{gathered}
F(z)={}_2F_1\left(\frac12,\frac12;1;z\right),\qquad
w(\delta)=\frac{(e^{\delta/2}-1)^3(e^{\delta/2}+3)}{16e^{\delta/2}},\\
G(\delta)=e^{3\delta/4}F(w(\delta))=\sum_{j\ge0}d_j\delta^j.
\end{gathered}
\tag{3}
\]

\textbf{Theorem 1 (all fixed late orders).} For each fixed positive
integer \(r\),

\[
\boxed{\displaystyle
 a_n=C\kappa^n\Gamma(n)
 \left\{\sum_{j=0}^{r-1}
 \frac{(-1)^j j!d_jL^j}{(n-1)_{\underline j}}+O_r(n^{-r})\right\}.}
\tag{4}
\]

Here \((n-1)_{\underline j}=(n-1)(n-2)\cdots(n-j)\), with the empty
product equal to one. In particular, (2) is true. The first terms are

\[
\begin{aligned}
\frac{a_n}{C\kappa^n\Gamma(n)}={}&1-\frac{3L}{4(n-1)}
+\frac{9L^2}{16(n-1)(n-2)}-\frac{15L^3}{32(n-1)(n-2)(n-3)}\\
&+\frac{135L^4}{256(n-1)\cdots(n-4)}
-\frac{903L^5}{1024(n-1)\cdots(n-5)}\\
&+\frac{63L^6}{32(n-1)\cdots(n-6)}+O(n^{-7}).
\end{aligned}
\tag{5}
\]

The proof below also establishes the exact Borel transform, an exact
median-Laplace representation, and a nonzero Stokes jump with action
\(L/4=\log3/2\) in the variable \(1/(4N)\). Thus the exponentially small
statement is not inferred merely from the divergent formal series.

\section{The moment integral and the exact local Borel
transform}\label{the-moment-integral-and-the-exact-local-borel-transform}

For independent uniform \(\theta,\phi\) on \([0,2\pi]\), put

\[
T=|1+e^{i\theta}+e^{i\phi}|^2\in[0,9].
\]

Expanding the powers and integrating the Fourier monomials gives
\(M_N=\mathbb E[T^N]\). Multiplying the three unit vectors by a common
independent phase does not change their resultant length. Consequently,
if \(p_3(x)\) is the radial density of the three-step planar uniform
random walk,

\[
M_N=\int_0^3x^{2N}p_3(x)\,dx,
\qquad \rho(t)=\frac{p_3(\sqrt t)}{2\sqrt t}
\tag{6}
\]

is the density of \(T\).

The classical elliptic formula, in the modulus convention
\(K(k)=\int_0^{\pi/2}(1-k^2\sin^2\theta)^{-1/2}\,d\theta\), is

\[
p_3(x)=\frac{\sqrt x}{\pi^2}\operatorname{Re}K\!\left(\sqrt{m(x)}\right),
\quad m(x)=\frac{(x+1)^3(3-x)}{16x}.
\tag{7}
\]

For \(1<x<3\), \(m(x)\in(0,1)\), so no real-part operation is needed. A
primary reference is Borwein--Straub--Wan--Zudilin, \emph{Densities of
Short Uniform Random Walks}, Canadian Journal of Mathematics 64 (2012),
961--990, equation (1.2). Its section 3, equation (3.1), explicitly
records the logarithmic singularity at \(x=1\).

For completeness, the part of (7) needed for the dominant endpoint can
be obtained directly. The two-step resultant has density
\(2/(\pi\sqrt{4-r^2})\). Conditional on length \(r\), adding a uniform
unit vector makes its squared resultant arcsine-distributed on
\([(r-1)^2,(r+1)^2]\). For \(1<x<3\), substitution \(s=r^2\) therefore
gives

\[
\rho(x^2)=\frac1{\pi^2}\int_{(x-1)^2}^{4}
\frac{ds}{\sqrt{s(4-s)((x+1)^2-s)(s-(x-1)^2)}}
=\frac{K(\sqrt{m(x)})}{2\pi^2\sqrt x},
\]

where the last equality is the standard four-real-root substitution for
the complete elliptic integral. Multiplication by \(2x\) gives (7) on
that interval.

Set \(t=9e^{-u}\), \(x=3e^{-u/2}\). On \(0\le u<L\), the normalized
density becomes

\[
f(u)=\frac{t\rho(t)}{C_0}
=e^{-3u/4}F\!\left(m(3e^{-u/2})\right),\qquad f(0)=1.
\tag{8}
\]

For the real density on all \(u\ge0\), use
\(f_{\rm phys}(u)=t\rho(t)/C_0\). Then

\[
Q(N):=\frac{NM_N}{C_0 9^N}
=N\int_0^\infty e^{-Nu}f_{\rm phys}(u)\,du.
\tag{9}
\]

The part away from \(u=0\) is exponentially small for any fixed positive
cutoff. Watson's lemma applied to the analytic germ (8) yields

\[
\boxed{a_n=4^n n![u^n]f(u).}
\tag{10}
\]

Thus the order-one Borel transform of \(\sum a_n z^n\), under the
convention \(\mathcal B\sum a_nz^n=\sum a_n\zeta^n/n!\), is exactly
\(f(4\zeta)\). Its nearest positive singularity will be \(\zeta=L/4\).

\section{Global continuation sufficient for coefficient
transfer}\label{global-continuation-sufficient-for-coefficient-transfer}

This step is essential: a real endpoint expansion by itself does not
prove (2).

The hypergeometric equation has possible singular values
\(m=0,1,\infty\). Its chosen solution \(F(m)\) is regular at \(m=0\).
The algebraic identities

\[
m(x)=\frac{-x^3+6x+8+3/x}{16},
\qquad 1-m(x)=\frac{(x-1)^3(x+3)}{16x}
\tag{11}
\]

give the following finite preimages under \(x=3e^{-u/2}\):

\begin{itemize}
\tightlist
\item
  \(m=1\): \(u=L+4\pi ik\), or \(u=(4k+2)\pi i\), \(k\in\mathbb Z\)
\item
  \(m=0\): \(u=4\pi ik\), or \(u=L+(4k+2)\pi i\)
\item
  \(m=\infty\): no finite \(u\), because the exponential never vanishes
\end{itemize}

Fix any \(R\) with \(L<R<2\pi\). In \(|u|<R\), the only singular-value
preimages are \(u=0\), where the initial branch is analytic, and
\(u=L\). Hypergeometric analytic continuation and the monodromy theorem
therefore continue (8) as a single-valued holomorphic function on

\[
\{u:|u|<R\}\setminus[L,R).
\tag{12}
\]

To be precise about the removable initial point: the germ at zero is
given by a convergent power series in \(m(u)\); any sufficiently small
loop around zero acts trivially on this germ. Continuation along paths
avoiding zero and \(L\), followed by the monodromy theorem in the slit
disk, is therefore compatible with its original holomorphic extension
through zero. Critical points of \(m(u)\) that do not map to a singular
value are merely pullback critical points, not singularities of the
composite function.

In particular \(u=L\) is the unique singularity on the circle of
convergence and the function is analytic in a standard dented disk (a
\(\Delta\)-domain) around it. No assumption of Borel continuation beyond
its justified domain has been made.

\section{The logarithmic multiplier and coefficient
transfer}\label{the-logarithmic-multiplier-and-coefficient-transfer}

The convergent zero-balanced hypergeometric expansion is

\[
F(1-w)=\frac1\pi\sum_{k\ge0}h_kw^k
\left\{2\psi(k+1)-2\psi(k+\tfrac12)-\log w\right\},
\qquad h_k=\frac{(1/2)_k^2}{(k!)^2}.
\tag{13}
\]

This is the \(a=b=1/2,m=0\) case of DLMF 15.8.10. Put \(\delta=L-u\).
Since \(x=e^{\delta/2}\), (11) yields

\[
1-m=w(\delta)=\frac{\delta^3}{32}(1+O(\delta)).
\tag{14}
\]

The factor \(w(\delta)/\delta^3\) is analytic and nonzero near zero.
With the branch continued from \(\delta>0\),

\[
\log w(\delta)=3\log\delta+\log\!\left(w(\delta)/\delta^3\right).
\]

Combining (8), (13), and \(e^{-3L/4}=1/(3\sqrt3)\), we obtain the exact
local form

\[
f(u)=-C G(L-u)\log(1-u/L)+H(u),
\tag{15}
\]

where \(H\) is holomorphic in a full neighborhood of \(L\). The
\(\log L\) term has been absorbed in \(H\). Equivalently, the analytic
logarithmic multiplier is

\[
g(u)=\frac3\pi e^{-3u/4}F(1-m(3e^{-u/2}))=C G(L-u),\quad g(L)=C.
\tag{16}
\]

For fixed integer \(j\ge0\) and \(n>j\), exact coefficient extraction
gives

\[
[u^n]\{(1-u/L)^j\log(1-u/L)\}
=\frac{(-1)^{j+1}j!L^{-n}}{n(n-1)\cdots(n-j)}.
\tag{17}
\]

Expand \(G(L-u)=\sum d_jL^j(1-u/L)^j\) and use (17). For a fully
explicit remainder proof, choose \(R_0\) with \(L<R_0<R\) sufficiently
close to \(L\) that the convergent local form (15) applies along the cut
from \(L\) to \(R_0\). A keyhole deformation of Cauchy's coefficient
contour, using the jump \(f_+-f_-=2\pi i C G(L-u)\), gives

\[
[u^n]f(u)=C\int_0^{R_0-L}G(-v)(L+v)^{-n-1}\,dv+O(R_0^{-n}).
\tag{17a}
\]

The small circle around \(L\) vanishes because the singularity is
logarithmic. On the outer contour the two boundary branches are bounded;
their contributions are \(O(R_0^{-n})\). Taylor's theorem gives
\(G(-v)=\sum_{j<r}(-1)^jd_jv^j+O_r(v^r)\). The remainder integral is
bounded by the full beta integral, and for \(n>j\),

\[
\int_0^\infty v^j(L+v)^{-n-1}\,dv
=L^{j-n}\frac{j!\,\Gamma(n-j)}{\Gamma(n+1)}.
\tag{17b}
\]

Extending each of the finitely many monomial integrals from
\([0,R_0-L]\) to \([0,\infty)\) changes it by an exponentially small
quantity \(O_j(R_0^{-n})\). Equations (17a)--(17b) prove

\[
[u^n]f(u)=\frac{CL^{-n}}n
\left\{\sum_{j<r}\frac{(-1)^jj!d_jL^j}{(n-1)_{\underline j}}+O_r(n^{-r})\right\}.
\]

There is no hidden logarithmic loss in this bound. Multiplication by
\(4^nn!\) proves Theorem 1.

Expanding (3) gives

\[
(d_0,\ldots,d_8)=\left(1,\frac34,\frac9{32},\frac5{64},\frac{45}{2048},
\frac{301}{40960},\frac7{2560},\frac{181}{172032},\frac{24251}{58720256}\right).
\tag{18}
\]

These numbers are exact rationals. Equation (3) is an all-order,
finite-at-each-order algorithm for all late correction coefficients.

\section{An exact Stokes sector and exact median
summation}\label{an-exact-stokes-sector-and-exact-median-summation}

Let \(f_+(u)\) and \(f_-(u)\), for real \(u>L\), denote analytic
continuation of the germ (8) past \(L\) in the upper and lower half
\(u\)-planes, respectively, with a small semicircle and then
continuation along the positive real axis. They are conjugates.

On the upper lip, \(\arg\delta=-\pi\). The analytic continuation of
\(\log w\) therefore has imaginary part \(-3\pi\), not \(-\pi\).
Equation (13) gives locally

\[
f_+(u)-f_-(u)=6i e^{-3u/4}F(1-m(3e^{-u/2}))=2\pi i g(u).
\tag{19}
\]

The right side continues without singularity along real \(u>L\), since
\(m>1\) and \(1-m<0\). Both sides solve the same pulled-back
second-order differential equation, so the local identity continues
along the entire interval. More concretely, if \(F(m+i0)\) denotes the
principal hypergeometric upper boundary value, then near \(L\), and
hence on the entire interval by continuation,

\[
f_+(u)=e^{-3u/4}\{F(m+i0)+2iF(1-m)\},
\quad f_-(u)=\overline{f_+(u)}.
\tag{20}
\]

Equivalently, standard hypergeometric continuation gives the especially
transparent formulas

\[
f_\pm(u)=e^{-3u/4}\{m^{-1/2}F(1/m)\ \pm\ 3iF(1-m)\},\qquad u>L.
\tag{20a}
\]

One way to verify the imaginary term is Pfaff's transformation
\(F(1-m)=m^{-1/2}F(1-1/m)\). The extra \(2iF(1-m)\) in (20) is purely
imaginary there. Equation (7) now proves the exact physical-median
identity

\[
f_{\rm phys}(u)=\frac{f_+(u)+f_-(u)}2\qquad(u>L).
\tag{21}
\]

It already holds trivially before \(L\), where both continuations equal
\(f\).

As \(u\to+\infty\), \(x=3e^{-u/2}\to0\) and \(m\asymp e^{u/2}\). The
standard coincident-exponent hypergeometric estimate gives
\(F(m\pm i0),F(1-m)=O(m^{-1/2}\log m)\). Thus

\[
f_\pm(u),\ g(u)=O((1+u)e^{-u}).
\tag{22}
\]

At \(u=L\), the singularity is logarithmic and integrable. The lateral
integrals

\[
Q_\pm(N)=N\int_0^\infty e^{-Nu}f_\pm(u)\,du
\tag{23}
\]

are therefore absolutely convergent for real \(N>0\), with the obvious
definition of \(f_\pm\) before \(L\). Equations (9) and (21) give

\[
\boxed{Q(N)=\frac{Q_+(N)+Q_-(N)}2.}
\tag{24}
\]

These are actual lateral lip (equivalently, indented-contour)
Borel--Laplace sums of (1), not only formal symbols: the Borel transform
is (10), continuation on the lateral ray is established by (20), and
growth is controlled by (22).

From (19), writing \(u=L+v\),

\[
\boxed{\displaystyle
Q_+(N)-Q_-(N)=\frac{2i}{\sqrt3}\,9^{-N}
N\int_0^\infty e^{-Nv}G(-v)\,dv.}
\tag{25}
\]

Consequently, for every fixed integer \(r\ge1\), Watson's lemma gives
the rigorous exponentially small expansion

\[
Q_+(N)-Q_-(N)=\frac{2i}{\sqrt3}\,9^{-N}
\left\{\sum_{j=0}^{r-1}\frac{(-1)^jj!d_j}{N^j}+O_r(N^{-r})\right\}.
\tag{26}
\]

In particular,

\[
Q_+-Q_-\sim\frac{2i}{\sqrt3}9^{-N}
\left(1-\frac3{4N}+\frac9{16N^2}-\frac{15}{32N^3}
+\frac{135}{256N^4}-\frac{903}{1024N^5}+\cdots\right).
\tag{27}
\]

This proves an exact nonzero Stokes sector with multiplier
\(2i/\sqrt3\). In the small parameter \(z=1/(4N)\), the exponential is
\(e^{-A/z}\) with \(A=L/4=\log3/2\). The same \(d_j\) control the late
terms in (4) and the Stokes-sector fluctuations in (26), an explicit,
rigorously derived resurgence relation.

\textbf{Scope boundary.} Equations (25)--(27) concern the difference of
lateral sums and their exact physical median. They do not assert that
the real moment function contains an independently adjustable real
\(9^{-N}\) term. Nor do they claim a complete global transseries,
numerical Stokes constants at the farther complex singularities
\(u=\pm2\pi i\), or uniform validity when the number of retained terms
grows with \(N\). Those questions require additional analysis.

\section{Inverting the late-coefficient
growth}\label{inverting-the-late-coefficient-growth}

Equation (4) implies \(a_n>0\) and \(a_{n+1}/a_n\sim\kappa n\); the
sequence is therefore strictly increasing eventually. The inverse
statement here is first stated on exact values \(y=a_n\), avoiding an
unjustified choice of continuous interpolation.

Rewrite (4) as

\[
a_n=C\kappa^n\Gamma(n)S(n),\qquad
S(n)\sim1+\frac{s_1}{n}+\frac{s_2}{n^2}+\cdots,
\]

where, for example,

\[
s_1=-\frac{3L}{4},\qquad
s_2=-\frac{3L}{4}+\frac{9L^2}{16}.
\tag{28}
\]

Stirling's expansion yields

\[
Y:=\log\frac{y}{C\sqrt{2\pi}}
=n\log\frac{\kappa n}{e}-\frac12\log n
+\sum_{j=1}^{K}\frac{c_j}{n^j}+O_K(n^{-K-1}),
\tag{29}
\]

with \(c_j\) computed by adding the formal logarithm of \(S\) to
Stirling's logarithmic correction series. In particular,

\[
c_1=\frac1{12}-\frac{3L}{4},\qquad
c_2=-\frac{3L}{4}+\frac{9L^2}{32}.
\tag{30}
\]

Let \(W=W_0\) be the positive real Lambert function and define

\[
t_0=\frac{Y}{W(\kappa Y/e)},\qquad h=\log(\kappa t_0),\qquad b=\frac{\log t_0}{2h}.
\tag{31}
\]

Then, as \(y=a_n\to\infty\),

\[
\boxed{\displaystyle n=t_0+b+
\frac{b/2-b^2/2-c_1}{t_0h}
+O\!\left(\frac1{t_0^2h}\right).}
\tag{32}
\]

Indeed, \(t_0\log(\kappa t_0/e)=Y\). Substitute \(t=t_0+b+\varepsilon\)
in (29): the leading residual is \((b^2/2-b/2+c_1)/t_0\), and the
derivative is \(h+O(t_0^{-1})\). The displayed \(\varepsilon\) cancels
that residual; the remaining residual is \(O(t_0^{-2})\), proving (32)
by the mean-value theorem. The derivative comparison uses an explicit
smooth asymptotic model, not an assumed interpolation of \(a_n\).

\textbf{All fixed inverse orders.} For arbitrary fixed \(K\), let
\(t_K\) be the unique sufficiently large solution of

\[
t\log(\kappa t/e)-\tfrac12\log t+\sum_{j=1}^{K}c_jt^{-j}=Y.
\tag{33}
\]

Its derivative is asymptotic to \(\log t\), so the same mean-value
argument gives

\[
n=t_K+O_K\!\left(\frac{n^{-K-1}}{\log n}\right)\quad(y=a_n).
\tag{34}
\]

The coefficients are effectively generated by (3), falling-factorial
expansion, logarithm, and Stirling's series. Newton iteration starting
at (31) produces the corresponding explicit nested-Lambert expansion to
any fixed order. For sufficiently large exact input \(a_n\), rounding
(32), or a sufficiently accurate \(t_K\), recovers \(n\).

\section{Inverting the original moment
sequence}\label{inverting-the-original-moment-sequence}

There is a natural continuous interpolation here:
\(M(s)=\mathbb E[T^s]\), real \(s>0\). It is eventually strictly
increasing, since the positive contribution from any interval
\(T\ge1+\varepsilon\) grows exponentially and dominates the bounded
negative contribution to \(M'(s)\) from \(0<T<1\). Differentiated Watson
estimates also show that the logarithmic derivative tends to \(L\).

For large \(y\), define

\[
n_0=-\frac1L W_{-1}\!\left(-\frac{C_0L}{y}\right),
\tag{35}
\]

so that \(C_0e^{Ln_0}/n_0=y\). Since

\[
\log\frac{NM_N}{C_0 9^N}=-\frac1{4N}+\frac1{32N^2}+O(N^{-3}),
\]

the actual inverse satisfies

\[
\boxed{\displaystyle M^{-1}(y)=n_0+\frac1{4Ln_0}
+\left(\frac1{4L^2}-\frac1{32L}\right)\frac1{n_0^2}
+O(n_0^{-3}).}
\tag{36}
\]

All fixed inverse orders follow by the same implicit-series/Newton
procedure from (1), with error divided by a derivative asymptotic to
\(L\). This is an inversion of the actual positive moment function.
Equations (35)--(36) do not by themselves claim a complete exponentially
improved inverse transseries.

\subsection{Exact inversion of the lateral sector
family}\label{exact-inversion-of-the-lateral-sector-family}

The exponentially small structure can also be inverted, rather than only
inverting its Poincaré part. Define

\[
J(s)=s\int_0^\infty e^{-sv}G(-v)\,dv,
\qquad D(s)=\frac{3}{4\pi}\int_0^\infty e^{-sv}G(-v)\,dv=\frac{3J(s)}{4\pi s},
\qquad \mathcal M_t(s)=M(s)+tD(s).
\tag{36a}
\]

The integral defines a holomorphic function for
\(\operatorname{Re}s>-1\), as does the moment integral \(M\). Equations
(24)--(25) imply exactly

\[
\mathcal M_i(s)=\frac{C_09^s}{s}Q_+(s),\qquad
\mathcal M_{-i}(s)=\frac{C_09^s}{s}Q_-(s).
\tag{36b}
\]

Let \(\phi=M^{-1}\), let \(s=\phi(y)\) for sufficiently large positive
\(y\), and choose the inverse branch \(n(t;y)\) of \(\mathcal M_t\) that
satisfies \(n(0;y)=s\). This branch exists and is analytic in \(t\) on a
disk of radius at least \(csy\), where \(c>0\) is independent of large
\(s\). Indeed, fix \(0<r<2\pi/L\). Uniform endpoint Laplace estimates
give

\[
\frac{M(s+w)}{M(s)}\longrightarrow e^{Lw}
\quad\text{uniformly for }|w|\le r,
\qquad D(s+w)=O(1/s).
\]

On \(|w|=r\), \(e^{Lw}-1\) has positive minimum modulus, and inside the
circle it has exactly one zero, counted with multiplicity. Rouché's
theorem first shows the same for \(M(s+w)-y\), and then, for
sufficiently small fixed \(c\), for \(M(s+w)+tD(s+w)-y\) uniformly on
\(|t|<csy\). The unique zero is simple, so the analytic implicit
function theorem gives the claimed analytic branch, with the uniform
bound \(|n(t;y)-s|<r\). This also proves the bounds needed for the
Cauchy remainder estimate below.

Lagrange--Bürmann inversion gives the exact convergent sector generator

\[
\boxed{\displaystyle
n(t;y)=\phi(y)+\sum_{k=1}^\infty\frac{(-t)^k}{k!}
\left(\frac d{dy}\right)^{k-1}
\left[\phi'(y)D(\phi(y))^k\right].}
\tag{36c}
\]

For example, apply ordinary reversion to \(z+tD(\phi(z))=y\), then the
generalized Lagrange formula to \(\phi(z)\). Its germ at \(t=0\) is
(36c), and analytic continuation on the established disk proves equality
throughout the disk of convergence. For each fixed \(K\), truncation at
\(k=K\), evaluated at either \(t=i\) or \(t=-i\), has rigorous error

\[
O_K((sy)^{-K-1}).
\tag{36d}
\]

This follows directly from Cauchy's estimate on a fixed fraction of the
disk \(|t|<csy\); no growing-order asymptotic assertion is being used.

Each fixed \(k\)-sector has an all-fixed-order expansion in \(s^{-1}\),
effectively obtained from

\[
\phi'(y)=\frac1{M'(s)},\qquad \frac d{dy}=\frac1{M'(s)}\frac d{ds},
\qquad J(s)\sim\sum_{j\ge0}(-1)^jj!d_js^{-j}
\]

and (1). The required differentiated remainder estimates follow by
differentiating the absolutely convergent Laplace integrals, or by
Cauchy estimates uniform on fixed disks centered at large positive
\(s\), retaining sufficiently many terms before taking a fixed number of
derivatives. Its leading term is

\[
\frac{(-t)^k}{k!}\partial_y^{k-1}[\phi'(y)D(\phi(y))^k]
=-\frac{t^k}{kL}\left(\frac{3}{4\pi sy}\right)^k(1+O_k(s^{-1})).
\tag{36e}
\]

Since \(sy=C_09^sQ(s)\), these are genuine \(9^{-ks}\) sectors. Thus
(36c)--(36e), together with the all-order fluctuation algorithm,
constitute a rigorous inversion of the exact lateral transseries family.

Writing \(n_\pm(y)=n(\pm i;y)\), the first terms are

\[
n_\pm(y)=s\mp i\frac{D(s)}{M'(s)}+O((sy)^{-2}),
\qquad n_+(y)-n_-(y)=-2i\frac{D(s)}{M'(s)}+O((sy)^{-3}).
\tag{36f}
\]

Odd/even cancellation in the second statement is justified by the
convergent analytic \(t\)-series, not merely by informal conjugation.
The first inverse Stokes sector therefore has the all-fixed-order
expansion beginning

\[
\boxed{\displaystyle
n_+(y)-n_-(y)= -\frac{3i}{2\pi Lsy}
\left[1+\frac{L^{-1}-3/4}{s}
+\frac{L^{-2}-L^{-1}+9/16}{s^2}+O(s^{-3})\right]
+O((sy)^{-3}).}
\tag{36g}
\]

The physical inverse is \(s=\phi(y)\). Inversion does not commute with
taking the median: from the even terms of (36c),

\[
\frac{n_+(y)+n_-(y)}2
=s-\frac12\frac d{dy}\left[\phi'(y)D(\phi(y))^2\right]+O((sy)^{-4}),
\tag{36h}
\]

whose leading correction is \(9/(32\pi^2L s^2y^2)\). This explicitly
distinguishes the physical inverse from the mean of the two lateral
inverses. The exact sector generator concerns this specified lateral
family; it does not assert that all possible distant-sheet sectors of
the original Borel transform have been classified.

% ed. (2026-10-02): names the repository apparatus that the three inversions re-derive without citing it (batch-77 placement dossier).
\begin{ednote}
The three inversions (31)--(34), (35)--(36) and (36c) are cases of the
inversion apparatus of ProveIt's canonical transseries volume,
\emph{Transseries: the polynomial-logarithmic calculus, series reversal
at infinity, and inversion} (Part~X, chapter \path{p0:sec:top}), which
the article does not cite (its repository searches in
Section~\ref{provenance-and-scope-of-the-novelty-check} looked for the
sequence numbers and ``honeycomb''); the method is the volume's, and only
the coefficients are new to the repository. The volume's source is
\path{Analysis/Transseries/docs/series-and-transseries/Transseries_And_Inversion/transseries_and_inversion.tex}.

\emph{Late coefficients.} The seed (31) solves the volume's factorial core,
Proposition~J.27 (\path{p0:prop:factorial-core}), with \(\kappa=1\),
\(d=\log\kappa-1\). In the displacement \(n=t_0(1+E)\), (33) becomes
\(hE+[(1+E)\log(1+E)-E]+f(t,E)=0\) with \(t=1/t_0\) and
\(f=-\tfrac12t\bigl(\log t_0+\log(1+E)\bigr)+\sum_{j\ge1}c_jt^{j+1}(1+E)^{-j}\),
which is its reversion about an arbitrary core, Theorem~J.33
(\path{p0:thm:core-reversion}), with \(\Lambda=h\); its first two
coefficients are \(b\) and \((b/2-b^2/2-c_1)/h\), the terms of (32), which
the editors recomputed symbolically. Recovering \(n\) by rounding when
\(y=a_n\) is part~(3) of the volume's staircase theorem, Theorem~J.43
(\path{p0:thm:staircase}): an approximation within \(\eta<\tfrac12\) of the
interpolated inverse at \(y=A_n\) rounds to \(n\).

\emph{Original moments.} Since \(\log M(s)=\log C_0+Ls-\log s+Q(1/s)\)
with \(Q(t)=-t/4+t^2/32+O(t^3)\), (35) is the volume's exact solution of
the dominant block, Theorem~J.23 (\path{p0:thm:lambert-core}), with
\(a=L\), \(b=-1\), the large root lying on the branch \(W_{-1}\) by its
branch rule (Corollary~J.25, \path{p0:cor:branch-rule}); and (36) is its
all-orders reversion around the Lambert core, Theorem~J.29
(\path{p0:thm:lambert-centered}), whose Bell recurrence gives
\(c_1=-q_1/a=1/(4L)\) and \(c_2=-(bc_1+q_2)/a=1/(4L^2)-1/(32L)\), the
coefficients of (36).

\emph{Sectors.} (36c) is the volume's perturbed inversion, Theorem~J.21
(\path{p0:thm:perturbed-inversion}), with \(F=M\), \(E=D\),
\(\varepsilon=t\) and \(\mu_0=\phi(y)\): its derivation
\(\mathcal D=F'(\mu)^{-1}\,d/d\mu\) is \(d/dy\), so its term of order
\(k\) (index \(M\) there),
\(\frac{(-t)^k}{k!}\mathcal D^{k-1}\bigl(D(\mu_0)^k/M'(\mu_0)\bigr)\), is
the \(k\)th term of (36c). The theorem is local; the radius \(csy\) of the
\(t\)-disk, and with it the remainder (36d), is the article's own Rouché
estimate.
\end{ednote}

\section{Exact replay and independent finite
checks}\label{exact-replay-and-independent-finite-checks}

The known recurrence for \(M_N\) is

\[
(N+1)^2M_{N+1}-(10N^2+10N+3)M_N+9N^2M_{N-1}=0.
\tag{37}
\]

Substitute the formal series (1), put \(B(z)=\sum_{n\ge0}a_nz^n/4^n\),
and compare coefficients in

\[
9(1+z)B\left(\frac z{1+z}\right)-(10+10z+3z^2)B(z)
+\frac1{1-z}B\left(\frac z{1-z}\right)=0.
\tag{38}
\]

This gives the exact triangular recurrence

\[
8na_n=\sum_{k=0}^{n-1}4^{n-k}
\left\{9\binom{1-k}{n+1-k}+\binom{n+1}{k}\right\}a_k-12a_{n-1},
\quad n\ge1,\quad a_0=1.
\tag{39}
\]

The binomial coefficient with negative top is the ordinary generalized
integer binomial. Uniqueness of the asymptotic coefficients and
recurrence matching justify (39). This is not used as a substitute for
the analytic proof of late growth.

\texttt{verify.py} uses Python integer/fraction arithmetic to:

\begin{enumerate}
\def\labelenumi{\arabic{enumi}.}
\tightlist
\item
  match all 22 displayed source terms exactly
\item
  compute the germ (8) independently of (39), matching 31 coefficients
  exactly
\item
  generate \(a_n\) through \(n=500\), checking integrality at every
  recurrence step
\item
  compute \(d_0,\ldots,d_{15}\) exactly
\item
  compare the late approximations against exact \(a_n\) at
  \(n=20,40,80,160,320,500\)
\end{enumerate}

The replay output is \texttt{verification.json}; the generated integers
are \texttt{coefficients\_0\_500.txt}. For example the exact-to-leading
ratio at \(n=500\) is approximately
\(0.996708446210591670313689669302\); retaining eight falling-factorial
terms leaves normalized error approximately \(2.45029\times10^{-18}\).
These finite numerical checks support implementation correctness; the
asymptotic theorems follow from sections 2--5.

\begin{center}
\begin{tabular}{rrrr}
\toprule
$n$ & $a_n/(C\kappa^n\Gamma(n))$ & Error after 4 terms & Error after 8 terms\\
\midrule
20 & 0.920460390900 & $1.08\cdot10^{-4}$ & $1.79\cdot10^{-6}$\\
40 & 0.959493066036 & $5.65\cdot10^{-6}$ & $2.92\cdot10^{-9}$\\
80 & 0.979570825175 & $3.25\cdot10^{-7}$ & $7.75\cdot10^{-12}$\\
160 & 0.989742589835 & $1.95\cdot10^{-8}$ & $2.52\cdot10^{-14}$\\
320 & 0.994860728387 & $1.20\cdot10^{-9}$ & $8.99\cdot10^{-17}$\\
500 & 0.996708446211 & $1.99\cdot10^{-10}$ & $2.45\cdot10^{-18}$\\
\bottomrule
\end{tabular}
\end{center}
The errors in this table are differences of normalized quantities, not relative errors of the original moments. Each column keeps a fixed number of terms.

The second replay, \texttt{verify\_analytic.py}, evaluates the density integral using a complementary-AGM form that preserves the cubic zero at $x=1$. At 70-digit working precision its relative errors against the exact moments at $N=1,2,5,10,20$ are below $10^{-50}$. These are floating-point consistency checks, not interval-certified error bounds. Formula (32) estimates the late index $n=500$ as $500.0000002664748\ldots$; formula (36) estimates the original moment index $N=30$ as $30.0000005839233\ldots$. Full numerical values are in \texttt{analytic\_verification.json}.

\section{Provenance and scope of the novelty
check}\label{provenance-and-scope-of-the-novelty-check}

Primary sources used:

\begin{itemize}
\tightlist
\item
  \href{https://oeis.org/A274600}{OEIS A274600}: target sequence,
  scaling, displayed terms, and explicit conjecture
\item
  \href{https://oeis.org/A002893}{OEIS A002893}: moment/count identity,
  honeycomb-return interpretation, recurrence, and older leading
  asymptotic
\item
  \href{https://people.mpim-bonn.mpg.de/zagier/files/doi/10.4153/CJM-2011-079-2/borweinJ0395.pdf}{Borwein--Straub--Wan--Zudilin
  (2012)}: equations (1.2), (3.1), and (3.2); elliptic density, its
  logarithmic singularity, and its relation to the same moments
\item
  \href{https://dlmf.nist.gov/15.8.E10}{NIST DLMF 15.8.10}: convergent
  zero-balanced hypergeometric logarithmic expansion
\item
  \href{https://dlmf.nist.gov/15.12}{NIST DLMF 15.12}: hypergeometric
  large-variable behavior
\item
  \href{https://cs.uwaterloo.ca/~shallit/Papers/cas.pdf}{Richmond--Shallit,
  Counting Abelian Squares}: the original abelian-square leading
  asymptotic is established background
\item
  \href{https://arxiv.org/abs/1609.05580}{Burnette--Wong, Abelian
  Squares and Their Progenies (2016)}: broader existing asymptotic
  context
\item
  \href{https://dirros.openscience.si/Dokument.php?id=41518&lang=eng}{Bašić--Fowler--Pickup--Potočnik
  (2026), Random walks and the electronic structure of graphene}:
  checked as a recent A002893 reference; its inverse-bandgap expansion
  uses the original moments, not the late coefficients in A274600
\item
  \href{https://math.stackexchange.com/questions/2006632/sum-involving-the-product-of-binomial-coefficients}{The
  Stack Exchange discussion linked by A274600}: 2016 discussion of the
  leading sum asymptotic; not a located proof of the late-coefficient
  assertion
\end{itemize}

Exact-ID searches for A274600 with proof, asymptotic, factorial,
ProveIt, and OEIS Open did not locate a proof. Searches for A002893 and
honeycomb/three-step late asymptotic coefficients did locate the
classical density literature and general resurgence literature, but no
source proving this exact late coefficient formula. The evidence
supports ``a proof of the explicitly labeled OEIS conjecture, with an
all-fixed-order and exact Stokes-jump strengthening''; it does not
establish that the mathematical mechanism was previously unknown or rule
out unpublished or differently indexed results. The elliptic density,
its van Hove logarithm, and Watson/singularity-transfer methods are
classical.

The source b-file was linked by the entry but inaccessible through the
checked HTTP/web paths; the finite source agreement is therefore
explicitly limited to the 22 displayed terms, not all 81 b-file entries.
The ProveIt default-branch code searches for A274600, A002893,
and honeycomb each returned zero results. The inspected complete
Analysis, Combinatorics, and Oeis subtree listings (5,850, 462, and 9
entries respectively; each marked untruncated) contain no exact
target-ID path. Their parent root tree was
4b874cea0012c51a6841ad9c58f6e20fa57c71da. This is a bounded duplicate
check, not a full-text search of every file at every historical
revision; details are in SOURCES.md. No external submission or
repository change has been made.

\section{Further questions}\label{further-questions}

The argument settles the labeled leading conjecture and gives all fixed
late orders, the exact positive-axis Stokes sector, and its
inverse-sector family. Several extensions remain outside its scope.

\begin{enumerate}
\def\labelenumi{\arabic{enumi}.}
\tightlist
\item
  Continue the selected Borel branch to the next singularities at
  \(u=\pm2\pi i\), determine their sheet-dependent connection constants,
  and isolate their effects on late coefficients after treating the
  dominant cut
\item
  Establish optimal truncation, quantitative exponentially improved
  remainder estimates, and uniform Stokes smoothing when the truncation
  order grows with the large parameter
\item
  Investigate all-index arithmetic properties of the normalized
  coefficients; the present replay checks integrality through index 500
  but does not supply a general integrality proof
\item
  Compare the same density-to-late-coefficient mechanism for larger
  alphabets, where different critical values and singularity types may
  govern the factorial scales
\end{enumerate}

These are research directions, not assertions that the corresponding
questions are new or unresolved in all existing literature.

\end{document}
